\subsection{Projected inversion and exclusion of asymmetric full matches}
The first label equation in \eqref{eq:independentmatch} is
\begin{equation}\label{eq:project}
 x(t,\eta)+y(t,\zeta)=0.
\end{equation}
For a possible source crossing, \(x(t,\zeta)\ge0\) and
\(y(t,\zeta)\le0\). The endpoint and monotonicity gates give
\[
 x(t,\zeta_a)+y(t,\zeta)>0,\qquad
 x(t,\zeta_b)+y(t,\zeta)<0,\qquad x_\eta<0.
\]
The intermediate value theorem and strict monotonicity therefore give one
and only one target phase \(\eta=\eta(t,\zeta)\) in the entire envelope.
In particular, eliminating this phase does not assume a branch near a
candidate point. The inverse is analytic wherever used, and
\begin{equation}\label{eq:inversepartial}
 \eta_\zeta=-\frac{y_\zeta(t,\zeta)}{x_\eta(t,\eta)},\qquad
 \eta_t=-\frac{y_t(t,\zeta)+x_t(t,\eta)}{x_\eta(t,\eta)}.
\end{equation}
On this projected graph, define the two remaining residuals
\begin{align}
 \HH(t,\zeta)&=x(t,\zeta)+y(t,\eta(t,\zeta)),\label{eq:HK}\\
 \KK(t,\zeta)&=q(t,\zeta)+q(t,\eta(t,\zeta))
                -\sigma(t)+d(x(t,\zeta),y(t,\zeta)).\notag
\end{align}
Projection together with \(\HH=\KK=0\) is equivalent to the complete
three-coordinate equation: if \(\HH=0\), both label pairs are reflected
and \(d(x(t,\eta),y(t,\eta))=d(x(t,\zeta),y(t,\zeta))\); then \(\KK=0\)
is exactly the complementary-\(q\) equation. Neither label equation alone
checks the amplitude ratio.

The contact identity \eqref{eq:contacttangent} and the projected equation,
without assuming \(\HH=0\), give
\begin{equation}\label{eq:HKderivatives}
 \HH_\zeta=x_\zeta(t,\zeta)
   -\frac{y_\eta(t,\eta)y_\zeta(t,\zeta)}{x_\eta(t,\eta)},\qquad
 \KK_\zeta=(y(t,\zeta)+1/2)\HH_\zeta .
\end{equation}
These exact identities control interval overestimation on projected cells.
Parameter derivatives instead use the full chain rule
\eqref{eq:inversepartial}; the fixed-parameter contact relation is not
applied to differentiation in \(t\).

\begin{certificate}[Exhaustive full-match exclusions]\label{cert:cover}
Let \(C_0=[0.0005,0.0007]\). Over the entire parameter box \(T^\Box\)
and all independent target phases in \([\zeta_a,\zeta_b]\), a closed
subdivision of the source envelope has 179 leaves. Each leaf passes one
of these predicates:
\begin{enumerate}
\item \(x<0\) or \(y>0\) throughout the cell (7 leaves);
\item the enclosure of \(\HH\) excludes zero (148 leaves);
\item the enclosure of \(\KK\) excludes zero (22 leaves);
\item the source cell lies in \(C_0\) and \(\HH_\zeta>0\)
throughout its projected family (2 leaves).
\end{enumerate}
A separate closed cover of \emph{all} of \(C_0\), with 554 leaves,
proves \(\HH_\zeta>0\) there. On \(C_0\), \(g=x+y\) has
\(g_\zeta<0\), with \(g(t,0.0005)>0\) and \(g(t,0.0007)<0\).
There are no unresolved leaves in either tree.
\end{certificate}

The finite input to each predicate is the analytic chart enclosure, not
sampled points. The inverse enclosure starts from the entire target
envelope and retains every solution under interval Newton intersection:
\[
 V_{\mathrm{new}}
 =V\cap\left(m-\frac{x(t,m)+y(t,\zeta)}{[x_\eta](t,V)}\right).
\]
The first derivative interval covers the full envelope. A floating-point
scout selects only a center, which is checked within the envelope; it does
not supply an exclusion or a target-phase restriction. Each retained
family has certified strict endpoint signs for the inverse. Subsequent
derivative intersections use already justified bounds. Tree acceptance
requires the children to have precisely the parent's endpoints and a
common join, every child to be reached exactly once, and every terminal
predicate to be resolved. These specifications are detailed in
Appendix~\ref{app:intervalspec}.

\begin{proposition}[Diagonal phases are a conclusion]\label{prop:diagonal}
For each \(t\in T\) there is a unique analytic phase
\(\phi(t)\in C_0\) with \(x(t,\phi)+y(t,\phi)=0\).
Every intrinsic full match has
\(\zeta=\eta=\phi(t)\).
\end{proposition}
\begin{proof}
The signs and nonzero phase derivative of \(g\) give its unique root in
\(C_0\), and the analytic implicit function theorem gives \(\phi\).
At that root \(\eta=\phi\) solves \eqref{eq:project}; uniqueness of
projected inversion identifies it as the target phase. Hence
\(\HH(t,\phi)=0\). The separate full central cover proves strict
monotonicity of \(\HH\) throughout \(C_0\), making \(\phi\) its only zero
there. The full cover excludes every noncentral full match by one of
the first three predicates. Lemma~\ref{lem:coverage} makes that cover
exhaustive for the intrinsic matching set, including both zero boundaries.
Thus every full match has the stated equal phases.
\end{proof}
Asymmetric solutions of the label equations are not discarded by
symmetry. The \(\KK\)-exclusion cells rigorously exclude the remaining
off-central label branches by their third-coordinate mismatch.

\subsection{Scalar equivalence and full-interval monotonicity}
We can now define
\begin{equation}\label{eq:residual}
 \Res(t)=2q(t,\phi(t))-\sigma(t)
           +d(x(t,\phi(t)),y(t,\phi(t))).
\end{equation}
Equations \eqref{eq:germ}, \eqref{eq:W}, the unique zero defining
\(\phi\), and \eqref{eq:residual} are an exact structural definition.
Truncated polynomials serve only to enclose this analytic function.

\begin{proposition}[Scalar/full-match equivalence]\label{prop:scalar}
For every \(t\in T\),
\[
 \Res(t)=0\quad\Longleftrightarrow\quad
 \mathcal M_{\sigma(t)}\ne\varnothing .
\]
\end{proposition}
\begin{proof}
Every full match is diagonal by Proposition~\ref{prop:diagonal}; its
third-coordinate equation is precisely \(\Res=0\).
Conversely, \(\Res(t)=0\) and \(g(t,\phi)=0\) give
\(x=-y\) and \(2q=\sigma-d(x,y)\). Thus \(W(t,\phi)=R_{\sigma(t)}W(t,\phi)\)
in all three coordinates. The positive-phase source tail and each of its
five propagated steps are valid and decreasing. Reflecting this chain
gives a valid decreasing forward tail to the negative fixed point.
Since \(x>y\) and \(x=-y\), the joined state has \(x>0>y\).
Applying \(\Phi_{\sigma(t)}^{-1}\) recovers every amplitude ratio and
places the join in the canonical \(\Sigma\), \(\mathcal U_s\) and
\(\mathcal V_s\). This is a genuine member of \(\mathcal M_s\).
\end{proof}

Implicit differentiation gives
\(\phi'=-(x_t+y_t)/(x_\zeta+y_\zeta)\).
Using \eqref{eq:contacttangent} in the full derivative of
\eqref{eq:residual} cancels the phase terms and yields
\begin{equation}\label{eq:residualderivative}
 \Res'(t)=2q_t-\sigma'(t)+(2x-1)y_t
 \quad\hbox{at }(t,\phi(t)).
\end{equation}
The partial derivatives on the right hold the phase fixed.

\begin{certificate}[Global derivative and endpoint signs]\label{cert:scalar}
On the entire closed box \(T^\Box\), the corrected analytic chart and
the interval enclosure of the whole phase family give
\[
 3.4800<\Res'(t)<3.5707.
\]
The analytic extension through \(a(L)\) satisfies
\[
 \Res(a(L))\in[-8.905044,-8.905043]\times10^{-16},\qquad
 \Res(a(U))\in[2.4652543,2.4652544]\times10^{-8}.
\]
These displayed outward bounds contain the full exact decimal intervals
in \pkg{scalar.json}; they are checked as rational comparisons.
\end{certificate}
This certificate encloses \(\phi(T^\Box)\) as an interval family and
evaluates \eqref{eq:residualderivative} on the full parameter interval.
It is not a set of derivative samples. Appendix~\ref{app:intervalspec}
gives the full derivative receipt and its finite acceptance rule.

\begin{theorem}[Global matching-parameter uniqueness]\label{thm:O6}
The projection
\[
 \pi_s(\mathcal M)=\{s\in I:\mathcal M_s\ne\varnothing\}
\]
is a singleton \(\{s_*\}\).
\end{theorem}
\begin{proof}
The strict derivative makes \(\Res\) increasing on \(T^\Box\), so it
has at most one zero there. Proposition~\ref{prop:scalar} and the
strict increase of \(\sigma\) then show that any two matching parameters
in \(I\) are equal. The strict endpoint signs and continuity give a zero
strictly between \(a(L)\) and \(a(U)\). The analytic extension through
the excluded left endpoint justifies this use of the intermediate value
theorem. Scalar equivalence supplies an intrinsic full match at that
zero. This existence calculation is independent of using
Proposition~\ref{prop:existmatch} as an existence premise, and uses no
fine root enclosure as an input.
\end{proof}
